\section{A dense elementary curve with rational witnesses}
\label{sec:density}

Define, for $1\le j\le n$,
\[
 h_j(t)=t\sin(t^{2j-1}).
\]
We prove quantitative density by maintaining rational intervals. The argument
uses only rational tests, finite grids, and explicit error bounds. In
particular, it never selects an exact zero of a real function.

The spacing of the exponents explains why one parameter can control all
coordinates. Around a large positive $T$, coordinate $j$ has amplitude
roughly $T$ and phase speed roughly $T^{2j-2}$. After coordinate $j-1$
has been controlled, the surviving interval has length comparable to
$\epsilon/T^{2j-3}$. Across that interval the next phase therefore changes
by an amount comparable to $\epsilon T$. Taking $T$ large creates room for
another positive and negative oscillation without destroying the earlier
coordinates. The bounds below turn this scale comparison into an explicit
rational construction; no equidistribution theorem is used.

\begin{theorem}[Primitive recursive rational density]\label{thm:density}
Let $a_1,\ldots,a_n\in\Q$, let $A\in\N$ satisfy
$\abs{a_i}\le A$, and let $0<\epsilon\le1$ be rational. There is a
primitive recursive construction of $q\in\Q$ such that
\[
 \abs{h_j(q)-a_j}<\epsilon/2\quad(1\le j\le n).
\]
One may choose an integer
\begin{equation}
 T=\left\lceil
       \max\{4(A+1),\,256(n+1),\,1024/\epsilon\}
    \right\rceil,
 \qquad B=T+64,\label{eq:density-size}
\end{equation}
and require $q\in[T,B]$.
\end{theorem}

\begin{proof}
We give the whole construction, including the finite searches.

\paragraph{A uniform Lipschitz bound.}
For $m=2j-1$ put
\[
 D_j=1+mB^m.
\]
The factorization of a difference of powers gives
$\abs{t^m-s^m}\le mB^{m-1}\abs{t-s}$ on $[T,B]$. Combining this with
$\abs{\sin v}\le1$ and
$\abs{\sin v-\sin w}\le\abs{v-w}$ yields
\begin{align}
 \abs{h_j(t)-h_j(s)}
 &\le\abs{t-s}+B\abs{t^m-s^m}\notag\\
 &\le D_j\abs{t-s}.\label{eq:density-lipschitz}
\end{align}
Thus no differentiation is needed to obtain the modulus used by the
construction.

Start with $I_0=[T,B]$. At stage $j$, we will construct a closed rational
interval $I_j\subseteq I_{j-1}$ with
\begin{equation}
 \abs{I_j}=\frac\epsilon{2D_j},
 \qquad
 \abs{h_i(t)-a_i}<\epsilon/2
 \quad(t\in I_j,\ i\le j).\label{eq:interval-invariant}
\end{equation}
Earlier coordinate estimates survive automatically when the interval is
shrunk. The main issue is to fit the next oscillation inside it.

\paragraph{Phase variation in the central half.}
Let $J=[u,v]$ be the central half of $I_{j-1}$. For $j=1$, this interval
has length 32, and the phase is $t$, so its phase variation is 32.
For $j\ge2$, again write $m=2j-1$. The previous exponent is $m-2$,
and factorization gives
\[
 v^m-u^m\ge mT^{m-1}(v-u)
 =\frac{mT^{m-1}\epsilon}{4D_{j-1}}.
\]
Since
$D_{j-1}=1+(m-2)B^{m-2}\le(m-1)B^{m-2}$, we obtain
\[
 v^m-u^m
 \ge\frac{\epsilon T}{4}
       \left(\frac TB\right)^{m-2}.
\]
For $0\le k\le2n-3$, Bernoulli's inequality and
$T\ge256(n+1)$ give
\[
 \left(\frac TB\right)^k
 =\left(1-\frac{64}{B}\right)^k
 \ge1-\frac{64k}{B}\ge\frac12.
\]
Consequently
\begin{equation}
 v^m-u^m\ge\epsilon T/8\ge128
 \quad(j\ge2).\label{eq:phase-room}
\end{equation}
The phase interval has width at least 32 in every stage.

\paragraph{Two rational points with opposite sine signs.}
Set $K=B^m+1$. Compute a rational $p$ satisfying
\[
 3<p<4,
 \qquad \abs{p-\pi}\le\frac1{64(K+1)}.
\]
The certified approximation primitive for $\pi$ supplies such a rational
with a precision bound computed from $K$. Find the least integer
$0\le k\le K$ for which
\[
 r_+=(2k+1/2)p\ge u^m+4.
\]
This is a bounded search with exact rational comparisons. Its last candidate
exceeds $u^m+4$, so it succeeds. The least candidate has $k\ge1$ because
$u\ge T$, and adjacent candidates differ by $2p<8$. Hence
\[
 u^m+4\le r_+<u^m+12.
\]
Put $r_-=(2k+3/2)p=r_++p$. Then $r_-<u^m+16$, so both rational
targets lie inside $[u^m,v^m]$, with substantial room to spare. Their
distances from the exact phases $(2k+1/2)\pi$ and $(2k+3/2)\pi$ are
less than $1/32$.

Subdivide $[u,v]$ into
\[
 N_0=1+\left\lceil16mB^{m-1}(v-u)\right\rceil
\]
equal subintervals. Consecutive rational phase values $t^m$ differ by at
most $1/16$. For each of $r_+$ and $r_-$, choose the first grid point
whose phase is at least that target. Both choices exist, since the final
phase is $v^m$. The chosen phase lies between its target and its target
plus $1/16$.

Call these points $q_+$ and $q_-$. Their order is $q_+<q_-$, since the
targets differ by $p>3$, much more than the possible overshoot. Their
phases differ from the corresponding exact sine maxima and minima by less
than $1/32+1/16<1/8$. Sine's Lipschitz bound now proves
\begin{equation}
 \sin(q_+^m)>7/8,
 \qquad
 \sin(q_-^m)<-7/8.\label{eq:opposite-sines}
\end{equation}
These inequalities were established from rational phase comparisons. They
were not obtained by invoking a procedure that decides an arbitrary real
sign.

\paragraph{Finding an approximately correct value.}
The size condition $T\ge4(A+1)$ makes
$h_j(q_+)-a_j$ positive by more than one and
$h_j(q_-)-a_j$ negative by more than one. Subdivide $[q_+,q_-]$ into
\[
 N_1=1+\left\lceil
          \frac{16D_j(q_--q_+)}\epsilon
       \right\rceil
\]
equal intervals. Write its grid points as $t_0,\ldots,t_{N_1}$.
At each one compute a rational $b_\ell$ with
\[
 \abs{b_\ell-(h_j(t_\ell)-a_j)}\le\epsilon/32.
\]
For example, approximate $\sin(t_\ell^m)$ to error
$\epsilon/(32B)$, multiply by the rational $t_\ell$, and subtract
$a_j$. This specifies a finite precision requirement.

The first $b_\ell$ is positive and the last is negative. Equation
\eqref{eq:density-lipschitz} and the grid spacing imply
\[
 \abs{b_\ell-b_{\ell-1}}
 \le\epsilon/16+2\epsilon/32=\epsilon/8.
\]
Choose the first index with $b_\ell\le0$. It has a predecessor with
positive value, so
\[
 -\epsilon/8<b_\ell\le0,
 \qquad
 \abs{h_j(t_\ell)-a_j}<5\epsilon/32<\epsilon/4.
\]
This is a finite discrete crossing argument. It constructs an approximate
solution at a rational point; an exact real solution is neither asserted
nor needed.

\paragraph{Restoring an interval.}
Set $c=t_\ell$ and define
\[
 I_j=\left[c-\frac\epsilon{4D_j},
           c+\frac\epsilon{4D_j}\right].
\]
Its length is the prescribed $\epsilon/(2D_j)$. For every $t\in I_j$,
the Lipschitz bound gives
\[
 \abs{h_j(t)-a_j}
 <5\epsilon/32+\epsilon/4
 =13\epsilon/32<\epsilon/2.
\]
The center $c$ lies in the central half of $I_{j-1}$. For $j\ge2$,
the elementary inequality $D_j\ge2D_{j-1}$ makes the new radius at most
one quarter of the length of $I_{j-1}$, so $I_j\subseteq I_{j-1}$.
For $j=1$, the radius is less than one and the available distance to the
endpoints is at least 16. This proves the same containment. The earlier
coordinates retain their estimates by nesting, establishing
\eqref{eq:interval-invariant}.

After $n$ stages, take the rational midpoint of $I_n$. It satisfies all the
claimed estimates and lies in $[T,B]$.
\end{proof}

\subsection{Why this is a primitive recursive construction}

All divisions, powers, floors, ceilings, and comparisons above are operations
on coded integers or rationals. Every search has a prescribed bound. For
example, $v-u\le64$ and $q_--q_+\le64$ give
\[
 N_0\le1024mB^{m-1}+2,
 \qquad
 N_1\le\frac{1024D_j}{\epsilon}+2.
\]
The second displayed upper bound may be replaced by its integer ceiling.
The approximations to $\pi$ and sine have primitive recursive truncation
bounds in their rational inputs and requested precision. A bounded loop over
$j=1,\ldots,n$, carrying the rational endpoints and its finite calculation
record, is therefore primitive recursive.

For an implementation, both grid traversals can be accelerated by binary
search on the grid indices. The phase values are monotone. In the second
search, monotonicity of $h_j$ is unnecessary: maintain endpoints whose
rational approximations have opposite signs, and keep the half with opposite
endpoint signs until the indices are adjacent. The same Lipschitz estimate
then certifies the desired approximation. The slower bounded traversal is
already sufficient for the proof.

Inside $\ipra$, the induction invariant is expressed by the finite record of
rational endpoint comparisons, approximation precisions, and rational error
bounds. One does not use induction on a formula quantifying over all real
points of the displayed intervals. The real notation in
\eqref{eq:interval-invariant} summarizes the consequence of these finite
certificates; the modulus \eqref{eq:density-lipschitz} is what permits this
summary to be unfolded at any requested rational precision.
